%!TEX root = ../../csuthesis_main.tex
\chapter{总结与展望}

\section{本文总结}

物理信息神经网络（Physics-Informed Neural Networks，PINN）是一种将深度学习技术与经典物理定律相结合的创新方法。这类神经网络不仅依赖于传统的数据驱动学习，还在其损失函数中融入了物理场的先验知识，从而能够在仅有少量观测数据甚至没有数据的情况下，解决各种复杂的物理问题，展现出了广泛的应用前景。然而，在实际应用中，PINN 模型的训练过程往往面临着收敛速度慢、计算效率低的问题。本文基于Navier-Stokes方程构建了PINN模型，并在此基础上对模型进行多方面的优化：
\begin{itemize}
    \item 在采样方法上，本文深入研究了基于残差的自适应采样方法。首先指出了传统固定和随机重采样技术在处理复杂PDE问题时的局限性。接着，深入介绍了RAR和RAD两种自适应策略，RAR根据残差大小动态添加采样点，RAD根据残差分布概率优化采样数据点，两者都致力于改善采样分布以精确匹配PDE解。后续针对RAD方法中的计算效率和参数调整难题，又提出了R3算法，有效地加快了模型收敛速度并提高了模型预测精度。
    \item 在损失函数上，本文针对PINN参数优化过程中的梯度不平衡问题进行改善。首先利用神经网络的自适应性将损失函数的动态权重加入训练，从而自适应地平衡各损失组成。随后，由加入一种新颖的残差注意力（RBA）机制，依靠累积残差的变化动态权衡采样点，无需额外梯度运算，进一步提升了求解效率与准确性。
    \item 在网络结构上，本文使用了一种改进的多层感知机架构，通过引入自注意力机制增强神经网络对输入特征的交互捕捉和深层表示能力，从而提升模型性能。
\end{itemize}

\section{未来展望}

本节就本文的研究工作以及PINN优化的近期研究现况， 提出以下值得进一步研究的问题:
\begin{itemize}
    \item 使用元学习改进采样策略：本文探讨了自适应采样在PINN训练过程中的重要性，RAD \cite{wu2023comprehensive}和R3 \cite{daw2022mitigating}是两种有效的方法，但两者都只在学习单个 PINN 时解决残差点级别的采样，Toloubidokhti等 \cite{toloubidokhti2024dats}提出的DATS则是实现了在元学习 PINN 时解决任务级采样，建立了第一个元学习 PINN 的自适应任务采样器，增强了模型对新任务或数据的适应性和泛化能力。
    \item 结合Transformer的网络结构：在本文中， PINN 始终依赖于多层感知器（MLP）进行训练，而忽略了实际物理系统中固有的关键时间依赖性，因此无法在全局范围内正确地传播初始条件约束并准确捕获各种场景下的真实解。Zhao等 \cite{zhao2023pinnsformer}基于 Transformer 模型中的多头注意机制捕获时间依赖性来准确地近似偏微分方程的解决方案，实现了更高的近似精度。
    \item 新的神经网络架构：Liu等 \cite{liu2024kan}受Kolmogorov-Arnold表示定理启发，提出一种MLP替代方案，称为Kolmogorov-Arnold Networks（KAN）。MLP将固定的激活函数放在节点（“神经元”）上，而KAN将可学习的激活函数放在边缘（“权重”）上，对于PDE求解，KAN比MLP更准确，收敛速度更快，损失更低。
\end{itemize}


